The four entries below organize the complete result by the action of \(J_1\)
and \(J_2\) on \(g\).  Equation~\eqref{eq:induced_time_reversals} fixes the
factor of \(\ii\) in each \(\Theta_a=U_a\mc K\).  Every displayed matrix also
obeys \(C\overline U_aC=U_a\).  Equations~\eqref{eq:group_only_generators} and
\eqref{eq:TR_canonical_and_invariant} then show that the two time reversals
leave the complete maximal subgroup invariant.

\cellentry{B04, B10, B12, F04, F10, F12}
{J_1^\trans gJ_1=g}
{J_2^\trans gJ_2=g}
{J_1J_2=-J_2J_1}
Here \(\Theta_1=J_1\mc K\) and \(\Theta_2=J_2\mc K\).  Physical
representations are
\begin{align*}
(J_1,J_2)_{\rm B04}={}&
 \bigl(I_2^{({\rm N})}\otimes\sigma_x^{({\rm f})},
 I_2^{({\rm N})}\otimes(\ii\sigma_y)^{({\rm f})}\bigr),\\
(J_1,J_2)_{\rm B10}={}&
 \bigl(I_2^{({\rm N})}\otimes(\ii\sigma_y)^{({\rm f})},
 I_2^{({\rm N})}\otimes\sigma_x^{({\rm f})}\bigr),\\
(J_1,J_2)_{\rm B12}={}&
 \bigl(I_2^{({\rm N})}\otimes\sigma_x^{({\rm f}_1)}
       \otimes(\ii\sigma_y)^{({\rm f}_2)},\notag\\[-2pt]
&\hspace{1.5em}
 I_2^{({\rm N})}\otimes(\ii\sigma_y)^{({\rm f}_1)}
       \otimes I_2^{({\rm f}_2)}\bigr),\\
(J_1,J_2)_{\rm F04}={}&
 \bigl(I_2^{({\rm N})}\otimes\sigma_x^{({\rm f})},
 I_2^{({\rm N})}\otimes(\ii\sigma_y)^{({\rm f})}\bigr),\\
(J_1,J_2)_{\rm F10}={}&
 \bigl(I_2^{({\rm N})}\otimes(\ii\sigma_y)^{({\rm f})},
 I_2^{({\rm N})}\otimes\sigma_x^{({\rm f})}\bigr),\\
(J_1,J_2)_{\rm F12}={}&
 \bigl(I_2^{({\rm N})}\otimes I_2^{({\rm f}_1)}
       \otimes(\ii\sigma_y)^{({\rm f}_2)},\notag\\[-2pt]
&\hspace{1.5em}
 I_2^{({\rm N})}\otimes(\ii\sigma_y)^{({\rm f}_1)}
       \otimes\sigma_x^{({\rm f}_2)}\bigr).
\end{align*}
The squares are \((+I,-I)\) for B04 and F04, \((-I,+I)\) for
B10 and F10, and \((-I,-I)\) for B12 and F12.  Thus B04, B10, F04,
and F10 realize the Majorana class, while B12 and F12 realize the Kramers
class.

\cellentry{B07, B13, B15, F07, F13, F15}
{J_1^\trans gJ_1=g}
{J_2^\trans gJ_2=-g}
{J_1J_2=J_2J_1}
Here \(\Theta_1=J_1\mc K\) and \(\Theta_2=\ii J_2\mc K\).  Complex
conjugation changes the sign of the explicit \(\ii\), so the commuting
matrices below induce anticommuting time reversals:
\begin{align*}
(J_1,J_2)_{\rm B07}={}&
 \bigl(I_2^{({\rm N})}\otimes I_2^{({\rm f})},
 \sigma_z^{({\rm N})}\otimes(\ii\sigma_y)^{({\rm f})}\bigr),\\
(J_1,J_2)_{\rm B13}={}&
 \bigl(I_2^{({\rm N})}\otimes(\ii\sigma_y)^{({\rm f})},
 \sigma_z^{({\rm N})}\otimes I_2^{({\rm f})}\bigr),\\
(J_1,J_2)_{\rm B15}={}&
 \bigl(I_2^{({\rm N})}\otimes(\ii\sigma_y)^{({\rm f})},
 \sigma_z^{({\rm N})}\otimes(\ii\sigma_y)^{({\rm f})}\bigr),\\
(J_1,J_2)_{\rm F07}={}&
 \bigl(I_2^{({\rm N})}\otimes I_2^{({\rm f})},
 (\ii\sigma_y)^{({\rm N})}\otimes I_2^{({\rm f})}\bigr),\\
(J_1,J_2)_{\rm F13}={}&
 \bigl(I_2^{({\rm N})}\otimes(\ii\sigma_y)^{({\rm f})},
 (\ii\sigma_y)^{({\rm N})}\otimes(\ii\sigma_y)^{({\rm f})}\bigr),\\
(J_1,J_2)_{\rm F15}={}&
 \bigl(I_2^{({\rm N})}\otimes(\ii\sigma_y)^{({\rm f})},
 (\ii\sigma_y)^{({\rm N})}\otimes I_2^{({\rm f})}\bigr).
\end{align*}
The squares are \((+I,-I)\) for B07 and F07, \((-I,+I)\) for
B13 and F13, and \((-I,-I)\) for B15 and F15.  Thus B07, B13, F07,
and F13 are in the Majorana class, while B15 and F15 are in the Kramers
class.

\cellentry{B19, B25, B27, F19, F25, F27}
{J_1^\trans gJ_1=-g}
{J_2^\trans gJ_2=g}
{J_1J_2=J_2J_1}
Now \(\Theta_1=\ii J_1\mc K\) and \(\Theta_2=J_2\mc K\).  The
representations are
\begin{align*}
(J_1,J_2)_{\rm B19}={}&
 \bigl(\sigma_z^{({\rm N})}\otimes I_2^{({\rm f})},
 I_2^{({\rm N})}\otimes(\ii\sigma_y)^{({\rm f})}\bigr),\\
(J_1,J_2)_{\rm B25}={}&
 \bigl(\sigma_z^{({\rm N})}\otimes(\ii\sigma_y)^{({\rm f})},
 I_2^{({\rm N})}\otimes I_2^{({\rm f})}\bigr),\\
(J_1,J_2)_{\rm B27}={}&
 \bigl(\sigma_z^{({\rm N})}\otimes(\ii\sigma_y)^{({\rm f})},
 I_2^{({\rm N})}\otimes(\ii\sigma_y)^{({\rm f})}\bigr),\\
(J_1,J_2)_{\rm F19}={}&
 \bigl((\ii\sigma_y)^{({\rm N})}\otimes(\ii\sigma_y)^{({\rm f})},
 I_2^{({\rm N})}\otimes(\ii\sigma_y)^{({\rm f})}\bigr),\\
(J_1,J_2)_{\rm F25}={}&
 \bigl((\ii\sigma_y)^{({\rm N})}\otimes I_2^{({\rm f})},
 I_2^{({\rm N})}\otimes I_2^{({\rm f})}\bigr),\\
(J_1,J_2)_{\rm F27}={}&
 \bigl((\ii\sigma_y)^{({\rm N})}\otimes I_2^{({\rm f})},
 I_2^{({\rm N})}\otimes(\ii\sigma_y)^{({\rm f})}\bigr).
\end{align*}
The squares are \((+I,-I)\) for B19 and F19, \((-I,+I)\) for
B25 and F25, and \((-I,-I)\) for B27 and F27.  Thus B19, B25, F19,
and F25 are in the Majorana class, while B27 and F27 are in the Kramers
class.

\cellentry{B24, B30, B32, F24, F30, F32}
{J_1^\trans gJ_1=-g}
{J_2^\trans gJ_2=-g}
{J_1J_2=-J_2J_1}
Here \(\Theta_1=\ii J_1\mc K\) and \(\Theta_2=\ii J_2\mc K\).  The
representations are
\begin{align*}
(J_1,J_2)_{\rm B24}={}&
 \bigl(\sigma_z^{({\rm N})}\otimes\sigma_x^{({\rm f})},
 \sigma_z^{({\rm N})}\otimes(\ii\sigma_y)^{({\rm f})}\bigr),\\
(J_1,J_2)_{\rm B30}={}&
 \bigl(\sigma_z^{({\rm N})}\otimes(\ii\sigma_y)^{({\rm f})},
 \sigma_z^{({\rm N})}\otimes\sigma_x^{({\rm f})}\bigr),\\
(J_1,J_2)_{\rm B32}={}&
 \bigl(\sigma_z^{({\rm N})}\otimes\sigma_x^{({\rm f}_1)}
       \otimes(\ii\sigma_y)^{({\rm f}_2)},\notag\\[-2pt]
&\hspace{1.5em}
 \sigma_z^{({\rm N})}\otimes(\ii\sigma_y)^{({\rm f}_1)}
       \otimes I_2^{({\rm f}_2)}\bigr),\\
(J_1,J_2)_{\rm F24}={}&
 \bigl((\ii\sigma_y)^{({\rm N})}\otimes(\ii\sigma_y)^{({\rm f})},
 (\ii\sigma_y)^{({\rm N})}\otimes\sigma_x^{({\rm f})}\bigr),\\
(J_1,J_2)_{\rm F30}={}&
 \bigl((\ii\sigma_y)^{({\rm N})}\otimes I_2^{({\rm f})},
 \sigma_z^{({\rm N})}\otimes I_2^{({\rm f})}\bigr),\\
(J_1,J_2)_{\rm F32}={}&
 \bigl((\ii\sigma_y)^{({\rm N})}\otimes I_2^{({\rm f})},
 \sigma_z^{({\rm N})}\otimes(\ii\sigma_y)^{({\rm f})}\bigr).
\end{align*}
The squares are \((+I,-I)\) for B24 and F24, \((-I,+I)\) for
B30 and F30, and \((-I,-I)\) for B32 and F32.  Thus B24, B30, F24,
and F30 are in the Majorana class, while B32 and F32 are in the Kramers
class.
